% FILE: 00_project_identity.tex
% Project: otel-weaver
% Role: telemetry-to-process-evidence-ingestion-layer

\section{Project Identity}
\label{sec:otel-weaver-identity}

\subsection{Purpose}
\label{sec:otel-weaver-identity-purpose}

The \texttt{otel-weaver} project is the telemetry-to-process-evidence ingestion layer of the
Process Intelligence Research Foundry. Its singular institutional purpose is to enforce a formal,
non-negotiable boundary between raw OpenTelemetry observability data---designated \emph{feedstock}
within the foundry---and court-grade process evidence, designated \emph{receipts}. This boundary is
not cosmetic. The project's foundational Feedstock Theorem states:
\begin{quote}
$\text{Telemetry} \in \text{Feedstock};\quad \text{SchemaValidation}(\text{Telemetry}) \neq
\text{ProcessCompliance}(\text{ProcessState})$
\end{quote}
The theorem is anchored at doctrine constant \texttt{TELEMETRY\_NOT\_EVIDENCE\_001} and appears
verbatim in \texttt{doctrine/telemetry-is-not-process-evidence.md}. The distinction it encodes is
operationally consequential: a passing Weaver schema check does not constitute a conformance
finding, and a conformance finding cannot be asserted without a signed receipt chain. The project
prevents \emph{category collapse}---the silent substitution of observability signal for process
truth---by requiring every ingestion path to traverse a typed admission gate before any downstream
process-evidence consumer may reference the payload.

The scope extends across five executable experiments, six immutable doctrine files, five YAML
mapping tables, a custom OTel \texttt{process.pi} semantic convention registry, and two
\texttt{ggen}-manufactured manifest chains. Each layer is designed to compose without gap: the
registry defines what feedstock looks like, the doctrine defines what feedstock is not, the mapping
tables define the translation rules, and the experiments materialize those rules as Rust type
systems and validated pipeline configurations.

\subsection{Repository Surface}
\label{sec:otel-weaver-identity-surface}

The project repository resides at \texttt{\$HOME/process-intelligence/otel-weaver/} and is
organized into the following primary directories:

\begin{description}
  \item[\texttt{doctrine/}] Six immutable doctrine files, each anchored at a named constant.
    Doctrine files are append-only by foundry policy and may not be rebased.
  \item[\texttt{experiments/}] Five subdirectories (\texttt{exp-001} through \texttt{exp-005}),
    each containing a \texttt{README.md} with full implementation specification, and in the cases
    of \texttt{exp-002} and \texttt{exp-003}, a \texttt{Cargo.toml} with real dependency
    declarations including \texttt{serde}, \texttt{blake3}, and a path reference to
    \texttt{wasm4pm-compat}.
  \item[\texttt{mappings/}] Five YAML mapping files governing OTel-to-PI witness translation,
    signal-to-evidence classification, registry-diff-to-drift disambiguation, live-check findings
    to refusal codes, and schema URL to receipt context.
  \item[\texttt{checkpoints/}] Three \texttt{ggen}-signed checkpoint files:
    \texttt{GGEN\_OTEL\_WEAVER\_PI\_PARTIAL\_001.md},
    \texttt{GGEN\_OTEL\_WEAVER\_PI\_ALIVE\_001.md}, and
    \texttt{GGEN\_OTEL\_WEAVER\_PI\_RUNTIME\_001.md}.
  \item[\texttt{ggen/}] Two manifest files and one audit script, governing the \texttt{ggen}
    manufacturing chain for this project.
  \item[\texttt{intel/}] Two YAML files mapping Weaver command surface and registry model
    structure.
\end{description}

The repository has no \texttt{.ttl}, \texttt{.rq}, or ontology files. The type surface is
expressed entirely through Rust type systems, YAML mapping tables, and doctrine constants rather
than through an RDF/SPARQL layer.

\subsection{Primary Research Role}
\label{sec:otel-weaver-identity-role}

Within the Process Intelligence Research Foundry, \texttt{otel-weaver} occupies the ingestion
stratum: it is the first component a raw OTel span touches before any process-intelligence
inference is permitted. No downstream component---including \texttt{wasm4pm},
\texttt{wasm4pm-compat}, or the M\&A claim layer---may consume telemetry-derived data unless it
has passed through an \texttt{Admission<T,\,W>} gate manufactured by this project's type
definitions.

This role is architecturally load-bearing. The Chatman foundry equation $A = \mu(O^*)$ requires
that every asserted activity $A$ be the manufactured output $\mu$ of an authoritative object state
$O^*$. The \texttt{otel-weaver} ingestion layer is the mechanism by which raw span data is either
promoted to a typed activity witness or refused with a structured \texttt{Refusal} diagnostic. No
middle ground is permitted by the type system.

\subsection{Key Artifacts}
\label{sec:otel-weaver-identity-artifacts}

The key artifacts manufactured or specified within this project are:

\begin{enumerate}
  \item \textbf{Admission\texttt{<T,\,W>}} --- A type-safe gate struct wrapping a verified
    feedstock payload \texttt{T} with a witness \texttt{W}. Defined in
    \texttt{exp-003-live-check-to-refusal/src/validator.rs}.
  \item \textbf{Refusal} --- An immutable structured diagnostic containing \texttt{RefusalCode},
    \texttt{violated\_rule}, \texttt{raw\_feedstock}, and a \texttt{blake3\_digest}.
  \item \textbf{BridgeRx} --- A Rust schema translation bridge mapping schema version diffs
    $\Delta W$ to a zero conformance residual $R_\Delta = 0$. Specified in
    \texttt{exp-002-weaver-diff-to-pi-residual/src/bridge.rs}.
  \item \textbf{ActivityWitness} --- A lattice-bounded Rust struct materialized from a resolved
    OTel Weaver registry JSON, anchoring the \texttt{OtelSpanToOcelEventProjection} witness type.
  \item \textbf{process.pi registry} --- A custom OTel semantic convention group defining
    \texttt{instance\_id}, \texttt{activity.name}, \texttt{lifecycle}, token state fields, and
    a \texttt{witness.hash} field constrained to BLAKE3 64-character digests.
  \item \textbf{LossReport} --- A mandatory declaration of lossy projection from OTel spans to
    OCEL events, covering \texttt{span\_duration\_truncation},
    \texttt{resource\_scope\_collapse}, and \texttt{causality\_to\_correlation\_degradation}.
\end{enumerate}
